\section{The two pure quadratic-eight characters at the quarter residues}
\label{sec:pure-quadratic8-quartic-strict}

\paragraph{Scope and inputs.}
The original finite characters here have trivial cubic part.
The two companion class statements of AC-G164-class, together
with the original quartic class inputs of AC-G158-class, give
exactly the class sectors used below.
The argument works directly on degree-four covers of
$K=\mathbb Q(\omega)$ and uses one explicit global unit for
each original quadratic character.
No conclusion for a trivial cubic part is inferred from the
earlier nontrivial cubic sectors.

Put $p=269$, and write
$\delta_s=\chi_{8s}|_{G_K}$ for $s=1,-1$.
Use the same symbols for the specified rational characters
when defining abelian fields.
Let $c_p$ be the actual cyclotomic character and $\omega_p$
its Teichmuller part.
Fix $i\in\mathbb Z_p$ with $i^2=-1$ and $\bar i=82$.
For a discrete coefficient $D$ use the full strict group
\[
 \begin{gathered}
 S_{\rm str}(K,D)=
 \ker\bigl(H^1_{\rm cts}(G_K,D)\longrightarrow\mathcal L(D)\bigr),\\
 \mathcal L(D)=
 \prod_{v\mid p}H^1(G_{K_v},D)
 \times\prod_{v\nmid p,\infty}H^1(I_v,D).
 \end{gathered}
 \tag{P8Q.1}
\]

\begin{theorem}[AC-G164: the two pure quadratic-eight quarter-residue rows]
\label{AC-G164}
Let $E/\mathbb Q_p$ be finite and let $O$ be its original
ring of integers, with arbitrary absolute ramification.
Let $s=1,-1$, $m\in\mathbb Z$, $e\geq1$, and let
$\kappa:G_K\to1+\mathfrak m_O$ be continuous and unramified
at every finite place away from $p$.
If $e+m\equiv66$ or $200\pmod{268}$, then
\[
 \begin{gathered}
 D=(E/O)(\delta_s\omega_p^{-m}\kappa^{-1}c_p^{-e}),\\
                       S_{\rm str}(K,D)=0.
 \end{gathered}
 \tag{P8Q.2}
\]
The residual strict groups vanish over every finite field
of characteristic $p$.
For $m=0,\kappa=1$, these are additional weights for the
two already present original finite labels $\chi_8,\chi_{-8}$,
both with trivial cubic part.
\end{theorem}

\begin{proof}
\proofstep{1. Actual covers and their selected characters.}
Put $\psi=\omega_p^{67}$, with $\psi(2)=-i$, and let
$\xi_s=\delta_s\psi$ as rational characters.
The exact quartic fields of AC-G158-class and AC-G164-class are
\[
 \begin{gathered}
 F_s=\mathbb Q(w_s),\qquad
        w_s^4+2spw_s^2+100p=0,\\
 F'_s=\mathbb Q(z_s),\qquad
        z_s=\sqrt{-3}\,w_s,\qquad
        z_s^4-6spz_s^2+900p=0,\\
                 N_s=KF_s=KF'_s.
 \end{gathered}
 \tag{P8Q.3}
\]
They are fields of $\xi_s$ and $\chi_{-3}\xi_s$,
respectively, with primitive conductors $8p$ and $24p$.
Both have unique quadratic subfield $\mathbb Q(\sqrt p)$,
different from $K$.
Thus $[N_s:\mathbb Q]=8$ and
$\Delta_s=\operatorname{Gal}(N_s/K)=C_4$.

Choose $b$ to fix $K$ and have $\xi_s(b)=-i$.
In the exact field coordinates,
\[
 \begin{aligned}
 b(w_s)&=-\frac{s w_s^3+438w_s}{130},\\
 b(z_s)&=\frac{s z_s^3-1314z_s}{390},\\
 b^2(w_s)&=-w_s,\qquad b^2(z_s)=-z_s.
 \end{aligned}
 \tag{P8Q.4}
\]
The CRT representative $5113$ is one modulo $24$ and two
modulo $p$.  It fixes $K$ and the quadratic conductor-eight
factors and supplies precisely this same orientation.
In particular these are automorphisms of the actual covers
over $K$, not involutions that conjugate the base.

Let $a$ denote the residual coefficient character of (P8Q.2),
and let $r\equiv e+m\pmod{268}$.
The Kummer quotient and selected inverse characters are
\[
 \begin{array}{c|c|c}
 r&\eta=a/\bar c_p&\eta^{-1}(b)\\ \hline
 66&\bar\xi_s^{-1}|_{G_K}&187=-\bar i\\
 200&\bar\xi_s|_{G_K}&82=\bar i .
 \end{array}
 \tag{P8Q.5}
\]
For example $-(66+1)=-67$ and $-(200+1)\equiv67\pmod{268}$;
the quadratic factors are their own inverses.
Thus $N_s/K$ splits $\eta$ with degree four, prime to $p$.
The quartic characters at $p$ occur only in these residual
splitting covers.  The original finite labels remain $\delta_s$.

\proofstep{2. Precisely the needed class sectors.}
The retained conclusions of AC-G158-class and AC-G164-class give
\[
       \operatorname{Cl}(F_s)[p]=
       \operatorname{Cl}(F'_s)[p]=0.
 \tag{P8Q.6}
\]
For clarity, the latter input uses an exact real quartic
unit-index relation $h_{F'_1}=2m$, $p\nmid m$, and an integral
Stickelberger annihilator for the odd companion $F'_{-1}$.
Neither input asserts a complete unit basis.

Let $c_K$ be the involution of $N_s$ fixing $F_s$ and
conjugating $K$.
It commutes with $b$, and $c_Kb^2$ fixes $F'_s$.
Norm to these two fixed fields followed by extension gives
\[
         1+c_K=0,\qquad1+c_Kb^2=0
                     \quad\hbox{on }\operatorname{Cl}(N_s)[p].
 \tag{P8Q.7}
\]
On the summand where $b^2=-1$, their sum is multiplication
by two.  Since $p$ is odd, it follows that
\[
                \operatorname{Cl}(N_s)[p]^{b^2=-1}=0.
 \tag{P8Q.8}
\]
No class-sector vanishing outside this summand is used.
Both inverse-character rows of (P8Q.5) lie in this summand.

\proofstep{3. Full-local Kummer descent and unit multiplicities.}
Put $k=O/\mathfrak m_O$.
Prime-to-$p$ restriction and inflation, including each local
decomposition and inertia quotient, identify the residual
strict group with
\[
 \begin{gathered}
 e_{\eta^{-1}}\ker\left(
 V_{N_s,S}\otimes_{\mathbb F_p}k
       \longrightarrow\mathcal P_{N_s}\otimes_{\mathbb F_p}k\right),\\
 V_{N_s,S}=\{x\in N_s^\times:
      p\mid v(x)\text{ for }v\nmid p\}/N_s^{\times p},\\
 \mathcal P_{N_s}=\bigoplus_{v\mid p}
                          N_{s,v}^{\times}/N_{s,v}^{\times p}.
 \end{gathered}
 \tag{P8Q.9}
\]
Over $N_s$ the coefficient is $k(1)$, with the extra quotient
twist $\eta$, so invariants select $\eta^{-1}$ in Kummer
cohomology.  There is no need to adjoin $\mu_p$.
The full local power condition also makes valuations above
$p$ divisible by $p$.  Thus the kernel lies in the all-primes
valuation group, with exact sequence
\[
 0\longrightarrow\mathcal O_{N_s}^{\times}/\mathcal O_{N_s}^{\times p}
 \longrightarrow V_{N_s}\longrightarrow
                          \operatorname{Cl}(N_s)[p]\longrightarrow0.
 \tag{P8Q.10}
\]
The selected class parts vanish by (P8Q.8); only units remain.

The four archimedean places above a fixed complex embedding
of $K$ form a regular $\Delta_s$-set.
The equivariant Dirichlet unit character is
$\mathbb Q[\Delta_s]-\mathbb Q$.
There are no roots of unity of order $p$ in $N_s$, since
$268\nmid[N_s:\mathbb Q]=8$.
Because $p\nmid4$, these multiplicities persist on the unit
quotient modulo $p$.
Consequently each of the $b$-eigenvalues $82$ and $187$
has a one-dimensional unit space over $\mathbb F_p$.
It suffices to give a nonzero full local Kummer image on
each of these four lines.

\proofstep{4. Two real fields and two exact global units.}
The totally real subfield used for $s=1$ is $F'_1$,
whereas for $s=-1$ it is $F_{-1}$.
Write $x_+=z_1$ and $x_-=w_{-1}$.
The exact equations and actions are
\[
 \begin{array}{c|c|c}
 &\text{equation}&b(x)\\ \hline
 +&x_+^4-1614x_+^2+242100=0&
                   (x_+^3-1314x_+)/390\\
 -&x_-^4-538x_-^2+26900=0&
                   (x_-^3-438x_-)/130 .
 \end{array}
 \tag{P8Q.11}
\]
Both equations are Eisenstein at $p$.
Their completions have ramification degree four.
Combining with the unramified quadratic $K_p$ gives the
full degree-eight completion of $N_s$.
Thus each $N_s$ has a unique place above $p$, of indices
$(e_p,f_p)=(4,2)$.
Each $x_s$ is a uniformizer, with $v_{x_s}(p)=4$,
and its four powers are an integral basis over
$\mathcal O_{K_p}$.

Use the following two global units:
\[
 \begin{aligned}
 U_+&=\frac{-292939446+7479912x_+
                         +1773428x_+^2-46473x_+^3}{78},\\
 U_-&=\frac{599770+81364x_-
                            -1150x_-^2-173x_-^3}{130}.
 \end{aligned}
 \tag{P8Q.12}
\]
They are unit two in the real companion certificate of
AC-G164-class and unit two in the $s=-1$ data of
AC-G158-class, respectively.
In the latter source every coefficient involving the cubic
generator is zero, so the element lies in $F_{-1}$ itself.
Their exact monic relations are
\[
 \begin{aligned}
 U_+^4-58370100U_+^3+5749535154002U_+^2
                                  +58370100U_++1&=0,\\
 U_-^4-8936U_-^3+5732922U_-^2+8936U_-+1&=0.
 \end{aligned}
 \tag{P8Q.13}
\]
The independent rational multiplication matrices in the basis
$1,x_s,x_s^2,x_s^3$ give exactly these characteristic
polynomials.
For $U_-$, the cube of its degree-four polynomial is also
checked to be the full degree-twelve characteristic polynomial
in the original source.
Monicity and constant coefficient one prove integrality of
the element and its inverse.
All rational denominators in (P8Q.12) are retained.
The verifier checks all three supplied real-field units in
each case; no fundamental-unit assertion is used here.

\proofstep{5. Four nonzero full local images.}
Reduction in $\mathbb F_{p^2}[x_s]/(x_s^4)$ gives
\[
 \begin{aligned}
 U_+&\equiv82+153x_++23x_+^2+56x_+^3,\\
 U_-&\equiv82+241x_-+136x_-^2+158x_-^3 .
 \end{aligned}
 \tag{P8Q.14}
\]
After removing the Teichmuller residue, direct finite
power-series logarithms through degree three are
\[
 \begin{aligned}
 \ell_+&=97x_++78x_+^3,\\
 \ell_-&=144x_-+211x_-^3 .
 \end{aligned}
 \tag{P8Q.15}
\]
For completeness, if $U/U(0)=1+y_1x+y_2x^2+y_3x^3$,
the three coefficients are
$y_1$, $y_2-y_1^2/2$ and
$y_3-y_1y_2+y_1^3/3$.
All denominators are invertible modulo $p$.
The complete nonlinear actions are
\[
       b(x_+)=187x_++249x_+^3,\qquad
       b(x_-)=187x_-+209x_-^3
                         \pmod{p,x_s^4}.
 \tag{P8Q.16}
\]
For $\lambda=82,187$, define
$P_\lambda=\sum_{j=0}^3\lambda^{-j}b^j$.
This is four times the corresponding idempotent on unit
classes modulo $p$.
Using the full actions in (P8Q.16) gives
\[
 \begin{aligned}
 P_{82}\ell_+&=245x_+^3,&
 P_{187}\ell_+&=119x_++67x_+^3,\\
 P_{82}\ell_-&=174x_-^3,&
 P_{187}\ell_-&=38x_-+132x_-^3 .
 \end{aligned}
 \tag{P8Q.17}
\]
The finite verifier checks all four values and their
$b$-eigenvalue identities.
Integer representatives of the coefficients of $P_\lambda$
give products of global conjugate units before localization.
Changing a representative by a multiple of $p$ changes
neither the Kummer class nor the displayed logarithm.

Since $4<p-1$, logarithm and exponential identify the local
principal units with the maximal ideal.
The logarithm of a principal $p$th power has valuation at
least five.
The four expressions in (P8Q.17) have nonzero leading
degree one or three, below that threshold.
Teichmuller units are themselves $p$th powers.
If a principal unit is a $p$th power in the full local
field, its root is a unit and its residue is one, since
Frobenius is bijective on the residue field.
Thus the four expressions are nonzero in the full local
power quotient, with no weakened local condition.

Localization is $\Delta_s$-equivariant.
Its nonzero image on each one-dimensional selected unit
space proves injectivity there, also after scalar extension
to any finite $k$.
Together with (P8Q.8)--(P8Q.10), this proves the residual
strict zero in both rows of (P8Q.5).

\proofstep{6. The original coefficient lift.}
Let $\varpi$ be a uniformizer of the original $O$.
The coefficient $D$ of (P8Q.2) is divisible and
$D[\varpi]\simeq k(a)$.
At $p$ the finite quadratic character is unramified and
$\kappa$ reduces to one.
Since $K_p/\mathbb Q_p$ is unramified, the residual inertia
characters $\bar c_p^{-66}$ and $\bar c_p^{-200}$ have
orders $134$ and $67$, respectively.
Both are nontrivial.
An inertia scalar differs from one by an $O$-unit, so
$D^{G_{K_p}}=0$.

At a finite place away from $p$, the cyclotomic,
Teichmuller and $\kappa$ factors are unramified.
The remaining quadratic inertia has order at most two,
prime to $p$.
Its invariants are zero or the whole divisible module $D$.
For every local group $J$ in (P8Q.1), therefore,
$D^J/\varpi D^J=0$.
The multiplication-by-$\varpi$ long exact sequence gives
local injectivity $H^1(J,D[\varpi])\to H^1(J,D)$.
A global strict class killed by $\varpi$ lifts to a residual
class, and local injectivity makes that lift strict.
It is zero by the residual assertion.
Every continuous cocycle from compact $G_K$ to discrete
$E/O$ has finite image and is killed by a power of $\varpi$.
A nonzero strict class would yield a nonzero one killed
by $\varpi$, which is impossible.
This proves (P8Q.2) over the original coefficient ring.
\end{proof}

\paragraph{Actual labels and exact added count.}
For $m=0,\kappa=1$ the finite labels are precisely
$\chi_8|_{G_K},\chi_{-8}|_{G_K}$, with trivial cubic part.
They occur on every prescribed host containing
$\mathbb Q(\zeta_8)$.
They remain distinct over $K$ since
$K\cap\mathbb Q(\zeta_8)=\mathbb Q$.
Their new weights have the form $e=66+134n$, $n\geq0$,
so their additional label-weight count below $T$ is
\[
                       2\left\lfloor\frac{T+67}{134}\right\rfloor.
 \tag{P8Q.18}
\]
Together with AC-G160, the union of their proved progressions
is $e\equiv-1\pmod{67}$, with count $2\lfloor T/67\rfloor$.
This statement concerns full strict groups.
Their possible positive derived indices and the actual
common-window transport still require the separate comparison.

\paragraph{Evidence and geometric scope.}
The exact finite script \path{verify_quartic_units.py} and
\path{G164_local_unit_certificate.json} record the six rational
unit relations, original source hashes, full local actions,
logarithms and projector values.
These are author finite checks for this theorem, separate
from the independent reviews and analytic replay of
AC-G164-class.
The latter theorem and AC-G158-class are used only for their
stated fields and class sectors.
The coefficients are original rank-one Artin--Tate systems
on the arithmetic curve of $K$.
The finite-cover dictionary remains SGA~1, Expos\'e~V,
\S\S4 and 7; Kummer theory retains the previously audited
Stacks tag~03PK scope.
The equivariant unit and class norm/extension inputs retain
their earlier exact source scopes.
No full normal-closure class group, complete unit basis,
mixed theta/Selmer realization, all-prime index bound,
H04.4 conclusion or Lean execution is asserted.
